\documentclass[11pt]{article}
\usepackage[margin=1in]{geometry}
\usepackage{hyperref}
\usepackage{graphicx}

% ── Paragraph spacing ────────────────────────────────────────────────────────
\usepackage{parskip}

% ── Page-break quality ───────────────────────────────────────────────────────
\widowpenalty=10000
\clubpenalty=10000
\raggedbottom

% ── Suppress automatic section numbering ─────────────────────────────────────
\setcounter{secnumdepth}{-1}

% ── Pandoc-generated table support ───────────────────────────────────────────
\usepackage{longtable}
\usepackage{booktabs}
\usepackage{array}
\usepackage{calc}
\usepackage{caption}
\newcounter{none}

% ── Math support ─────────────────────────────────────────────────────────────
\usepackage{amsmath}
\usepackage{amssymb}

% ── Make \paragraph and \subparagraph block-level ────────────────────────────
\usepackage{titlesec}
\titleformat{\paragraph}
  {\normalfont\normalsize\bfseries}{\theparagraph}{1em}{}
\titlespacing*{\paragraph}
  {0pt}{3.25ex plus 1ex minus .2ex}{1.5ex plus .2ex}
\titleformat{\subparagraph}
  {\normalfont\normalsize\bfseries}{\thesubparagraph}{1em}{}
\titlespacing*{\subparagraph}
  {0pt}{3.25ex plus 1ex minus .2ex}{1.5ex plus .2ex}

% ── Pandoc-generated tightlist support ───────────────────────────────────────
\providecommand{\tightlist}{%
  \setlength{\itemsep}{0pt}\setlength{\parskip}{0pt}}

% ── Pandoc 3.x figure sizing ────────────────────────────────────────────────
\newcommand{\pandocbounded}[1]{%
  \sbox0{#1}%
  \ifdim\wd0>\linewidth
    \resizebox{\linewidth}{!}{#1}%
  \else
    #1%
  \fi
}

% ── Unicode character fixes ──────────────────────────────────────────────────
\usepackage{newunicodechar}
\newunicodechar{✓}{$\checkmark$}
\newunicodechar{≠}{$\neq$}
\newunicodechar{≡}{$\equiv$}
\newunicodechar{≈}{$\approx$}
\newunicodechar{δ}{$\delta$}
\newunicodechar{−}{$-$}
\let\mystRightarrow\rightarrow
\renewcommand{\rightarrow}{\ensuremath{\mystRightarrow}}
% ─────────────────────────────────────────────────────────────────────────────

\hypersetup{colorlinks=true, allcolors=blue}

\title{Technical Summary: Value Dynamics}
\author{Keith Lostracco}
\date{}

\begin{document}
\maketitle

\subsection{Model}\label{model}

A \textbf{high-energy center} (HEC) is a network node with accumulated
value mass \(\mathcal{M} > 0\) (energy units). The \textbf{coupling
distance} \(r \in (0, \infty)\) between an agent and a HEC measures the
degree of the agent's independence (\(r \to 0\): maximal coupling;
\(r \to \infty\): minimal coupling). This is a generalized coordinate,
not a spatial distance.

The agent's \textbf{net energy rate} at coupling distance \(r\) is:

\[\Pi(r) = \underbrace{\frac{G\mathcal{M}}{r}}_{\text{resource inflow}} - \underbrace{\frac{\tau\mathcal{M}}{r^2}}_{\text{dissolution cost}} - \underbrace{\gamma B_i}_{\text{entropy maintenance}}\]

where \(G > 0\) is the resource coupling coefficient, \(\tau > 0\) is
the dissolution coupling coefficient, \(\gamma > 0\) is the entropy
leakage rate, and \(B_i > 0\) is boundary integrity. The
\textbf{coexistence potential} is \(V(r) = -\Pi(r)\):

\[V(r) = \frac{\tau\mathcal{M}}{r^2} - \frac{G\mathcal{M}}{r} + \gamma B_i\]

This potential has exactly one critical point on \((0, \infty)\), a
global minimum at the \textbf{coexistence attractor} \(r^* = 2\tau/G\),
with minimum value \(V(r^*) = \gamma B_i - G^2\mathcal{M}/(4\tau)\). The
well depth is \(D = G^2\mathcal{M}/(4\tau)\).

Agents adjust coupling via gradient dynamics \(\dot{r} = -\mu V'(r)\),
where \(\mu > 0\) is the adjustment rate. The boundary is an
\textbf{active dissipative structure} with reference integrity
\(\bar{B} > 0\), collapse floor \(B_{\min} \in (0, \bar{B})\) (the set
\(\{B_i \leq B_{\min}\}\) is absorbing: reaching it is dissolution), and
structural-capacity factor
\(\zeta(B_i) = \max((B_i - B_{\min})/(\bar{B} - B_{\min}), 0) \in [0, 1]\),
the fraction of repair and mobility capacity retained at integrity
\(B_i \in (B_{\min}, \bar{B}]\). Boundary dynamics couple to coupling
distance through assimilation pressure
\(D_{\text{assimilate}}(r) = \sigma\mathcal{M}/r^3\) and autocatalytic,
thermodynamically priced repair
\(R_{\text{repair}}(r, B_i) = (\rho/\kappa)\,\zeta(B_i)\,\max(\Pi(r), 0)\),
where \(\kappa > 1\) is the repair multiplier of TC-IV. The coupled
system, for \(B_i > B_{\min}\):

\[\dot{r} = -\mu_0\,\zeta(B_i)\,V'(r), \qquad \dot{B}_i = \frac{\rho}{\kappa}\,\zeta(B_i)\,\max(\Pi(r), 0) - \frac{\sigma\mathcal{M}}{r^3} - \gamma B_i\]

with baseline mobility \(\mu_0 > 0\); the constant-\(\mu\) dynamics are
recovered in the healthy regime with \(\mu = \mu_0\zeta(B_i^*)\).
\textbf{Proposition (Healthy Boundary Equilibrium):} at fixed \(r\)
satisfying the repair-viability condition
\(q\Pi_0(r) > \gamma + 2\sqrt{q\gamma d_0(r)}\) --- where
\(q = \rho/(\kappa(\bar{B} - B_{\min}))\),
\(\Pi_0(r) = G\mathcal{M}/r - \tau\mathcal{M}/r^2 - \gamma B_{\min}\),
\(d_0(r) = \sigma\mathcal{M}/r^3 + \gamma B_{\min}\) --- the boundary
dynamics have exactly two fixed points above the floor: an unstable
\(B_{\mathrm{u}}(r)\) and a stable \(B_i^*(r)\) (saturated at
\(\bar{B}\) if it exceeds the reference integrity), with \(\Pi > 0\) at
\(B_i^*(r)\). Under the viable-agent condition (repair-viability at
\(r^*\) with \(B_{\mathrm{u}}(r^*) < \bar{B}\)), the point
\((r^*, \min(B_i^*(r^*), \bar{B}))\) is a locally asymptotically stable
equilibrium of the coupled system, and the band/attractor results hold
verbatim at \(B_i = B_i^*\).

\subsection{Results}\label{results}

\textbf{Theorem 22 (Stability of the Coexistence Attractor).} Under
gradient dynamics \(\dot{r} = -\mu V'(r)\), the coexistence attractor
\(r^* = 2\tau/G\) is:

\begin{enumerate}
\def\labelenumi{(\alph{enumi})}
\item
  A fixed point: \(\dot{r}|_{r=r^*} = 0\).
\item
  Locally asymptotically stable: eigenvalue
  \(\lambda = -\mu G^4\mathcal{M}/(8\tau^3) < 0\).
\item
  Globally attracting on \((0, \infty)\): every trajectory starting at
  \(r_0 \in (0, \infty)\) converges to \(r^*\) as \(t \to \infty\).
\end{enumerate}

\textbf{Theorem 23 (Existence of the Coexistence Band).} The Stable
Coexistence Band \(\mathcal{B} = (r_-, r_+)\) exists if and only if:

\[\mathcal{M} > \mathcal{M}_{\min} \triangleq \frac{4\gamma B_i \tau}{G^2}\]

The band boundaries are
\(r_\pm = (G\mathcal{M} \pm \sqrt{\Delta})/(2\gamma B_i)\) where
\(\Delta = G^2\mathcal{M}^2 - 4\gamma B_i \tau \mathcal{M}\). When
\(\mathcal{M} = \mathcal{M}_{\min}\), the open band is empty:
\(V \geq 0\) everywhere with equality only at \(r^*\) (marginal,
non-surviving). When \(\mathcal{M} < \mathcal{M}_{\min}\), no viable
coupling distance exists.

\textbf{Theorem 24 (Freedom Bandwidth).} The width of the Stable
Coexistence Band is:

\[w \triangleq r_+ - r_- = \frac{\sqrt{\Delta}}{\gamma B_i} = \frac{\sqrt{G^2\mathcal{M}^2 - 4\gamma B_i \tau \mathcal{M}}}{\gamma B_i}\]

The bandwidth satisfies:

\begin{enumerate}
\def\labelenumi{(\alph{enumi})}
\item
  \(w > 0\) iff \(\mathcal{M} > \mathcal{M}_{\min}\).
\item
  \(\partial w / \partial \mathcal{M} > 0\) (richer centers offer wider
  freedom).
\item
  \(w\) is decreasing in \(\gamma\), \(B_i\), and \(\tau\).
\item
  \(w\) is increasing in \(G\).
\item
  For \(\mathcal{M} \gg \mathcal{M}_{\min}\):
  \(w \approx G\mathcal{M}/(\gamma B_i)\) (linear growth).
\item
  Near \(\mathcal{M}_{\min}\):
  \(w \sim \sqrt{\mathcal{M} - \mathcal{M}_{\min}}\) (square-root
  vanishing).
\end{enumerate}

\textbf{Corollary 24.1 (Freedom Is Finite).} For any finite
\(\mathcal{M}\), \(w\) is finite. Infinite freedom requires
\(\mathcal{M} \to \infty\) or \(\gamma B_i \to 0\) (physically
impossible).

\textbf{Theorem 25 (Basin of No Return).} Let
\(\bar{\Pi} \triangleq G^2\mathcal{M}/(4\tau) - \gamma B_{\min}\)
(surplus ceiling) and let \(B_c\) be the smaller root of
\(q(B_i - B_{\min})(G^2\mathcal{M}/(4\tau) - \gamma B_i) = \gamma B_i\).
If repair-viability holds at some \(r\), \(B_c\) exists with
\(B_{\min} < B_c \leq B_{\mathrm{u}}(r)\); for any agent admitting a
healthy equilibrium, moreover
\(B_c \leq B_{\mathrm{u}}(r^*) < \bar{B}\), and:

\begin{enumerate}
\def\labelenumi{(\alph{enumi})}
\item
  If \(B_i(t_0) < B_c\), then \(\dot{B}_i < 0\) along every trajectory
  regardless of the path \(r(t)\), and \(B_i\) reaches \(B_{\min}\) in
  finite time; no later surplus can reverse the decline.
\item
  As \(B_i \downarrow B_{\min}\), mobility dies:
  \(\mu(B_i) = \mu_0\zeta(B_i) \to 0\).
\item
  Above the basin, recovery is possible: from \(r(t_0) = r^*\) with
  \(B_i(t_0) > B_{\mathrm{u}}(r^*)\), the agent converges to the healthy
  equilibrium. \(B_c\) is a certified inner bound of the basin.
\end{enumerate}

\textbf{Theorem 26 (Irreversibility of Dissolution).} Define the
dissolution threshold
\(r_d \triangleq (\kappa\sigma\mathcal{M}/(\rho\bar{\Pi} + \kappa\gamma B_{\min}))^{1/3}\),
i.e.~\(\sigma\mathcal{M}/r_d^3 = (\rho/\kappa)\bar{\Pi} + \gamma B_{\min}\).
Then:

\begin{enumerate}
\def\labelenumi{(\alph{enumi})}
\item
  Whenever \(r(t) < r_d\): \(\dot{B}_i < -\gamma(B_i + B_{\min}) < 0\)
  for every \(B_i \in (B_{\min}, \bar{B}]\).
\item
  If \(r(t) < r_d\) for all \(t \geq t_0\) (in particular if
  \(r_d \geq r^*\), by (c)), then \(B_i\) reaches \(B_{\min}\) at a
  finite \(t_d \leq t_0 + (B_i(t_0) - B_{\min})/(2\gamma B_{\min})\):
  the agent dissolves.
\item
  If \(r_d \geq r^*\), every trajectory starting below \(r_d\) remains
  below \(r_d\) (the attractor lies inside the dissolution zone) and (b)
  applies unconditionally; if \(r_d < r^*\) and \(\ell > \ell_c\)
  (below), \(B_i\) falls below \(B_c\) before escape and collapse is
  irreversible by Theorem 25.
\end{enumerate}

\emph{Remark (Race-Condition Closure).} For \(r_d < r^*\), the
phase-space bound \(dB_i/dr \leq -\alpha\) with
\(\alpha = \sigma\kappa\gamma B_{\min}/(2\mu_0\tau(\rho\bar{\Pi} + \kappa\gamma B_{\min})) = \gamma B_{\min} r_d^3/(2\mu_0\tau\mathcal{M})\)
holds on \((0, r_d)\); dissolution is irreversible when the penetration
depth \(\ell = r_d - r_0\) exceeds
\(\ell_c = (B_i^{(0)} - B_c)/\alpha\): boundary depletion crosses the
basin before the gradient escape reaches \(r_d\). The bound is
conservative --- mobility death (\(\mu(B_i) \to 0\)) makes the true
decline per unit escape distance diverge as \(B_i \to B_{\min}\) --- so
\(\ell > \ell_c\) is sufficient, far from necessary.

\textbf{Corollary 26.1 (Assimilation Trap).} \(r_d\) and the energetic
inner boundary \(r_-\) are distinct and either ordering can occur. If
\(r_d > r_-\), an agent at \(r \in (r_-, r_d)\) earns positive net
energy (\(\Pi(r) > 0\)) while losing identity (\(\dot{B}_i < 0\)); as
\(B_i\) falls, \(r_-\) moves inward, so an agent that entered below
\(r_-\) becomes energetically viable while remaining inside the
dissolution zone --- irreversibly once \(B_i < B_c\).

\textbf{Theorem 27 (Starvation Spiral).} If \(r > r_+\) (outside the
Coexistence Band):

\begin{enumerate}
\def\labelenumi{(\alph{enumi})}
\item
  \(\Pi(r) < 0\): energy deficit, no repair.
\item
  \(\dot{B}_i = -\sigma\mathcal{M}/r^3 - \gamma B_i < 0\), with
  \(B_i(t) \leq B_i^{(0)}e^{-\gamma t}\) while the deficit lasts; in the
  isolation limit the collapse floor is reached at the finite time
  \(t^* = (1/\gamma)\ln(B_i^{(0)}/B_{\min})\).
\item
  Recovery channels --- inward drift (re-entry from \(r_0\) while
  \(B_i > B_c\) takes at least
  \((r_0^3 - r_+(B_c)^3)/(3\mu_0\mathcal{M}G)\), the expanding boundary
  remaining inside \(r_+(B_c)\)) and outward expansion of \(r_+(B_i)\)
  as the maintenance load falls --- both close at \(B_c\).
\item
  If \(B_i < B_c\) --- sufficient condition:
  \(r_0^3 \geq r_+(B_c)^3 + (3\mu_0\mathcal{M}G/\gamma)\ln(B_i^{(0)}/B_c)\)
  --- mobility and repair vanish together, the drift stalls before
  re-entry, and \(B_i\) reaches \(B_{\min}\) in finite time:
  irreversible starvation. Briefly starved agents above the basin can
  recover.
\end{enumerate}

\textbf{Corollary 27.1 (Inequality of Freedom).} If \(B_i > B_j\) (agent
\(i\) is more complex): (a) \(w_i < w_j\); (b)
\(\mathcal{M}_{\min,i} > \mathcal{M}_{\min,j}\); (c) for any
\(\mathcal{M}\) with
\(\mathcal{M}_{\min,j} < \mathcal{M} < \mathcal{M}_{\min,i}\), agent
\(j\) survives near the HEC but agent \(i\) cannot.

\textbf{Theorem 28 (Multi-Center Attractor).} The multi-center
coexistence potential
\(V(\mathbf{r}) = \sum_{k=1}^K [\tau_k\mathcal{M}_k/r_k^2 - G_k\mathcal{M}_k/r_k] + \gamma B_i\)
has a unique global minimum at \(r_k^* = 2\tau_k/G_k\) for each \(k\),
which is globally attracting under gradient dynamics. The multi-center
Coexistence Band exists iff:

\[\sum_{k=1}^{K} \frac{G_k^2 \mathcal{M}_k}{4\tau_k} > \gamma B_i\]

\textbf{Corollary 28.1 (Diversification Benefit).} An agent may be
viable when coupled to \(K\) HECs even if no single HEC has sufficient
mass
(\(\mathcal{M}_k < \mathcal{M}_{\min,k} \triangleq 4\gamma B_i \tau_k / G_k^2\)
individually), provided the combined well depths exceed \(\gamma B_i\).

\textbf{Corollary (Cascade Collapse).} If a HEC with mass
\(\mathcal{M}\) is destroyed, all \(N\) agents within its Coexistence
Band simultaneously lose viability with respect to that center (agents
coupled solely to it lose viability outright). Total freedom destroyed:
\(\mathcal{F}_{\text{lost}} = \sum_i w_i(\mathcal{M})\), reducing to
\(N \cdot w(\mathcal{M})\) for homogeneous agents.

\subsection{Notation}\label{notation}

{\def\LTcaptype{none} % do not increment counter
\begin{longtable}[]{@{}
  >{\raggedright\arraybackslash}p{(\linewidth - 2\tabcolsep) * \real{0.5000}}
  >{\raggedright\arraybackslash}p{(\linewidth - 2\tabcolsep) * \real{0.5000}}@{}}
\toprule\noalign{}
\begin{minipage}[b]{\linewidth}\raggedright
Symbol
\end{minipage} & \begin{minipage}[b]{\linewidth}\raggedright
Meaning
\end{minipage} \\
\midrule\noalign{}
\endhead
\bottomrule\noalign{}
\endlastfoot
\(\mathcal{M}\) & HEC value mass (accumulated negentropy, J) \\
\(r\) & Coupling distance (generalized coordinate) \\
\(G\) & Resource coupling coefficient \\
\(\tau\) & Dissolution coupling coefficient \\
\(\gamma\) & Entropy leakage rate \\
\(B_i\) & Agent boundary integrity (energy units) \\
\(V(r)\) & Coexistence potential:
\(\tau\mathcal{M}/r^2 - G\mathcal{M}/r + \gamma B_i\) \\
\(r^*\) & Coexistence attractor: \(2\tau/G\) \\
\(\mathcal{M}_{\min}\) & Minimum viable HEC mass:
\(4\gamma B_i \tau / G^2\) \\
\(r_-, r_+\) & Energetic inner boundary and outer (starvation) threshold
of the band \\
\(w\) & Freedom bandwidth: \(r_+ - r_-\) \\
\(D\) & Potential well depth: \(G^2\mathcal{M}/(4\tau)\) \\
\(\mu_0\), \(\mu(B_i)\) & Baseline mobility; health-dependent mobility
\(\mu_0\zeta(B_i)\) \\
\(\sigma\) & Assimilation intensity coefficient \\
\(\rho\) & Repair allocation fraction \\
\(\kappa\) & Repair multiplier (\(> 1\), TC-IV) \\
\(\bar{B}\), \(B_{\min}\) & Reference integrity; collapse floor
(absorbing below) \\
\(\zeta(B_i)\) & Structural-capacity factor:
\(\max((B_i - B_{\min})/(\bar{B} - B_{\min}), 0)\) \\
\(\bar{\Pi}\) & Surplus ceiling:
\(G^2\mathcal{M}/(4\tau) - \gamma B_{\min}\) \\
\(B_c\) & Basin boundary (basin of no return below) \\
\(r_d\) & Dissolution threshold:
\((\kappa\sigma\mathcal{M}/(\rho\bar{\Pi} + \kappa\gamma B_{\min}))^{1/3}\) \\
\(\ell_c\) & Critical penetration depth: \((B_i^{(0)} - B_c)/\alpha\) \\
\end{longtable}
}

\end{document}
